\section{具体的な可換模型と中心式の可換化の相違}
\label{em:sec:comparison}

可換復元には、候補となる方程式を取り出すことと、その方程式の
Kuznetsov成分へのenhanced同値を作ることの二段階がある。
一楕円ブロックの分離型では、前者は二次関係式の多重線形化から計算でき、
後者には定理\ref{er:thm:main}のブロック比較を使える。
本節ではこの具体性を明示し、三次中心元の表示から積の順序を消す操作が、
その可換模型を与えないことを証明する。
一般混合型への圏同値は本節の結論に含めない。
Hessianによる点曲線公式はMori--Smithの既知の結果であり、
本節では量子化の規約に合わせた証明を示し、既存の圏比較への接続を明示する。
以下の二次比較と混合項の障害も初等的に証明できる。
それらを先行研究に対して新規なものとする確認は未完了である。

\subsection{三変数対称Jacobian量子化のHessian公式}

三変数の任意の斉次三次式$f\in\CC[x,y,z]_3$に対し、
$\operatorname{Hess}(f)=\det(\partial_i\partial_j f)$と書く。
これはHessian\emph{行列}ではなく、その行列式である。
以下の多項式恒等式は$f$の滑らかさを仮定しない。
点曲線の記述はMori--Smith
\cite[Lemma 7.1, Theorem 7.2, pp.~23--25]{PMoriSmith}にある。
同論文の$A_{f,\lambda}$の規約では、行列式は
$\lambda^3\operatorname{Hess}(f)/8+3\lambda f$である。
ここでは関係式の符号を合わせて、その式を$-3h$で正規化する。

\begin{proposition}[Mori--SmithのHessian公式の量子化表示]
\label{em:prop:hessian-formula}
対称Jacobian関係式
\[
 [y,z]=h\sym(\partial_xf),\qquad
 [z,x]=h\sym(\partial_yf),\qquad
 [x,y]=h\sym(\partial_zf)
\]
を多重線形化すると、$v=(x,y,z)^{\mathsf T}$について係数行列は
\[
 M_f(v)=K(v)-\frac h2 H_f(v),\qquad
 K(v)=\begin{pmatrix}0&-z&y\\z&0&-x\\-y&x&0\end{pmatrix},
 \quad H_f=(\partial_i\partial_j f)
\]
となる。その行列式は
\begin{equation}\label{em:eq:hessian-identity}
 \det M_f=-3h\left(f+\frac{h^2}{24}\operatorname{Hess}(f)\right).
\end{equation}
従って$h\ne0$では、多重線形化の退化三次式は
\begin{equation}\label{em:eq:hessian-candidate}
 f_h^{\mathrm{pt}}=f+\frac{h^2}{24}\operatorname{Hess}(f)
\end{equation}
で与えられる。この式は、退化三次式が恒等的に零になる場合も含む恒等式である。
実際の点曲線との同定には、対応するAS-regular代数と点加群についての追加仮定が必要である。
\end{proposition}
\begin{proof}
二次式$q$の対称化の多重線形化は
$\frac12\sum_j(\partial_jq)(v)v'_j$である。
$q=\partial_i f$と置けば行列表示を得る。
任意の対称$3\times3$行列$H$と$\varepsilon$について、
$K^{\mathsf T}=-K$より
\[
 \det(K-\varepsilon H)=-\det(K+\varepsilon H).
\]
左辺は$\varepsilon$の三次以下の奇多項式であり、
$\operatorname{adj}K=vv^{\mathsf T}$なので一次係数は
$-\operatorname{tr}(\operatorname{adj}K\,H)=-v^{\mathsf T}Hv$である。
三次係数は$-\det H$であり、従って
\[
 \det(K-\varepsilon H)=-\varepsilon v^{\mathsf T}Hv-
 \varepsilon^3\det H.
\]
三次のEuler等式を二回用いると$v^{\mathsf T}H_fv=6f$である。
$\varepsilon=h/2$を代入すれば\eqref{em:eq:hessian-identity}を得る。
\end{proof}

この公式は座標選択とも整合する。実際、
$g(x)=\lambda f(Lx)$、$L\in\mathrm{GL}_3(\CC)$、$\lambda\ne0$なら
\begin{equation}\label{em:eq:hessian-covariance}
 \operatorname{Hess}(g)(x)
 =\lambda^3\det(L)^2\operatorname{Hess}(f)(Lx),\qquad
 g_h^{\mathrm{pt}}(x)=\lambda f_{h'}^{\mathrm{pt}}(Lx),
 \quad h'=\lambda\det(L)h.
\end{equation}
第一式は$H_g=\lambda L^{\mathsf T}H_f(Lx)L$の行列式を取れば従い、
第二式は代入で得られる。$h'$はJacobian Poisson括弧の座標変換則と一致する。
固定した体積形式の下での式であり、Hessianの数値係数を座標変換と独立に扱わない。

\begin{corollary}[一般の滑らかな一ブロックの候補式]
\label{em:cor:hessian-reconstruction}
$g$は滑らかな平面三次曲線を定め、$\Omega_h^g$は
系\ref{er:cor:ternary}の正規化された中心元とする。
同系のパラメータ開集合と$P(u)$の仮定の下では、
定理\ref{er:thm:main}の圏比較を入力とすることにより、
\begin{equation}\label{em:eq:general-hessian-output}
 V\left(g(x)+\frac{h^2}{24}\operatorname{Hess}(g)(x)+P(u)\right)
\end{equation}
が$S_h^g[u]/(\Omega_h^g+P)$のrotation付き可換模型となる。
これは同系の模型の射影同型な表示であり、一般混合型の復元定理ではない。
\end{corollary}
\begin{proof}
$g=\lambda f_t(Lx)$とHesse形に移し、$h'=\lambda\det(L)h$とする。
直接計算で
\[
 \operatorname{Hess}(f_t)=
 2(t^3+4)xyz-2t^2(x^3+y^3+z^3),\qquad
 f_t+\frac{h'^2}{24}\operatorname{Hess}(f_t)
 =b(h',t) f_{\vartheta(h',t)}
\]
を得る。従って\eqref{em:eq:hessian-covariance}により
\eqref{em:eq:general-hessian-output}の方程式は
$\lambda b(h',t) f_{\vartheta(h',t)}(Lx)+P(u)$となる。
$b(h',t)\ne0$なので、三変数を一斉に$b(h',t)^{-1/3}$倍すれば、
系\ref{er:cor:ternary}の方程式
$\lambda f_{\vartheta(h',t)}(Lx)+P(u)$になる。
この座標変換と同系の圏同値により結論を得る。
ここで用いた三乗根は複素ファイバーごとに選べる。
族上の同型を選ぶ際には、必要な有限\'etale被覆を別途指定する。
\end{proof}

\subsection{二次関係式から取り出せる分離型の方程式}

$S_{h,t}$、$\Omega_{h,t}$、$P(u)$および$\mathscr U_t$は
第\ref{er:sec:intro}節のものとする。
$t^3\ne-1$、$h\in\mathscr U_t$、$m\ge2$とし、$P$のアフィン臨界点は原点のみとする。
三次中心商は
\[
 B_{h,t,P}=S_{h,t}[u_1,\ldots,u_m]/(\Omega_{h,t}+P(u))
\]
である。既存のSklyanin中心商と楕円曲線の比較
\cite[pp.~20--21]{SmithTate}\cite[Sections 3.5--3.6]{EG}だけでなく、
内部シフトとの整合性および中心変数を加えた後の生成性を使う点が重要である。

\begin{proposition}[分離型の具体的復元手順]
\label{em:prop:determinant}
$\alpha=1-ht/2$、$\beta=-1-ht/2$、$\gamma=-h$と置く。
関係式を第二の三変数について線形にした行列を
\[
 M(x,y,z)=
 \begin{pmatrix}
 \gamma x&\beta z&\alpha y\\
 \alpha z&\gamma y&\beta x\\
 \beta y&\alpha x&\gamma z
 \end{pmatrix}
\]
とすると
\begin{equation}\label{em:eq:determinant}
 -\det M=\alpha\beta\gamma(x^3+y^3+z^3)
 -(\alpha^3+\beta^3+\gamma^3)xyz.
\end{equation}
$h\ne0$では、$b=1-h^2t^2/4$を用いて
\begin{equation}\label{em:eq:explicit-output}
 F_{\mathrm{rec}}(h)
 =-\frac{\det M}{3hb}+P(u)
 =\frac{x^3+y^3+z^3}{3}+\vartheta(h,t)xyz+P(u)
\end{equation}
を得る。右辺は$h=0$まで正則に延長する。
定理\ref{er:thm:main}の圏比較を入力とすれば、
$Y_h=V(F_{\mathrm{rec}}(h))$は$B_{h,t,P}$のrotation付き可換模型である。
\end{proposition}
\begin{proof}
多重線形化された関係式は
$M(x,y,z)(x',y',z')^{\mathsf T}=0$である。
行列式を六項に展開すると\eqref{em:eq:determinant}を得る。
さらに
\[
 \alpha\beta\gamma=hb,\qquad
 \alpha^3+\beta^3+\gamma^3
 =-3h\left(t+\frac{h^2}{3}(1+t^3/4)\right)
\]
を代入すれば\eqref{em:eq:explicit-output}となる。
零パラメータでは行列式全体が零になるため、共通因子$h$を除いてから延長する。
これは$S_{0,t}$のpoint schemeを楕円曲線と同一視する操作ではない。
最後の圏同値は定理\ref{er:thm:main}そのものであり、
行列式の計算だけから従うものではない。
\end{proof}

一般の三次元Sklyanin関係式
$\alpha yz+\beta zy+\gamma x^2=0$とその巡回置換で与えられた場合も、
滑らかな点曲線および標準中心三次商であることを確認したうえで、
\eqref{em:eq:determinant}が点曲線の方程式を直接与える。
分離型では方程式の非零スカラー倍はブロック変数の一斉拡大で吸収できる。
その拡大に必要な三乗根は$\CC$上で存在するが、一般の基礎体上では底拡大の指定が必要である。
これは元の古典極限の方程式を再利用する手順ではなく、
与えられた量子関係式から変形した楕円曲線を取り出す手順である。

\subsection{中心式の積を可換にすると二次で違う模型になる}

$S_3=x^3+y^3+z^3$と置き、中心元の指定された対称表示
$\Omega_{h,t}=U(h,t)S_3/3+W(h,t)\sym(xyz)$の積の順序を消した多項式を
\[
 F_{\mathrm{sym}}(h)=U(h,t)S_3/3+W(h,t)xyz+P(u)
\]
と定義する。これは選んだ表示に依存する操作であって、
代数の可換化とは異なる。可換化との相違は次の小節で扱う。

\begin{proposition}[二次Jacobian類の不一致]
\label{em:prop:second-order-failure}
$F_0=f_t+P$、$J(F_0)=\CC[x,y,z,u]/(\partial F_0)$とする。
$F_{\mathrm{rec}}(h)$と$F_{\mathrm{sym}}(h)$はいずれも$h$の一次項を持たず、
二次係数のJacobian類はそれぞれ
\begin{equation}\label{em:eq:two-classes}
 \frac{1+t^3}{3}[xyz],\qquad
 -\frac{2(1+t^3)}{3}[xyz]
 \quad\text{in }J(F_0)_3
\end{equation}
である。差$(1+t^3)[xyz]$は非零である。
従って二つの偏極付きcubic族は、$\CC[h]/(h^3)$上で
特殊ファイバーの恒等写像を保つ射影座標変換によって同型にならない。
特に、中心式の積の順序を消す操作は、分離型でもmarkedな可換復元を与えない。
\end{proposition}
\begin{proof}
式\eqref{er:eq:coefficients}を展開すると
\begin{align*}
 \vartheta(h,t)&=t+\frac{1+t^3}{3}h^2+O(h^4),\\
 U(h,t)&=1+\frac{t^2}{2}h^2+O(h^4),\\
 W(h,t)&=t-\frac{t^3+4}{6}h^2+O(h^4).
\end{align*}
$J(f_t)$では
$x^2+tyz=y^2+tzx=z^2+txy=0$なので
$S_3=-3txyz$である。従って$F_{\mathrm{sym}}$の二次係数の類は
\[
 \left[\frac{t^2}{6}S_3-\frac{t^3+4}{6}xyz\right]
 =-\frac{2(1+t^3)}{3}[xyz]
\]
となる。

$[xyz]\ne0$も有限な線形代数で直接確認できる。
次数三の単項式を
\[
 (x^3,x^2y,x^2z,xy^2,xyz,xz^2,y^3,y^2z,yz^2,z^3)
\]
の順に並べる。列を
$(x\partial_xf_t,y\partial_xf_t,z\partial_xf_t,
x\partial_yf_t,y\partial_yf_t,z\partial_yf_t,
x\partial_zf_t,y\partial_zf_t,z\partial_zf_t,xyz)$
とする係数行列の行列式は
\[
 -(1+t^3)^2.
\]
従って$t^3\ne-1$のとき、最初の九列は独立であり、$xyz$はその張る空間に含まれない。
変数の分離から$J(F_0)=J(f_t)\otimes J(P)$であり、
$J(P)_0=\CC$なので$[xyz]\otimes1$も非零である。

最後に、$F_0$は滑らかなcubicなので、その偏微分は二次の正則列をなす。
これらに一次係数のsyzygyはなく、Euler等式を除いて無限小射影自己同型はない。
従って一次項が零である二つのmarked族の間の一次座標変換は射影的に零である。
二次座標変換と方程式の単元倍で変えられる二次項は$(\partial F_0)_3$に限られる。
差のJacobian類が非零なので、求める同型は存在しない。
\end{proof}

\begin{example}[Fermatブロック]
\label{em:ex:fermat}
$t=0$、$P(u,v)=u^3+v^3$では
\begin{align*}
 F_{\mathrm{rec}}(h)&=\frac{x^3+y^3+z^3}{3}+u^3+v^3
                       +\frac{h^2}{3}xyz,\\
 F_{\mathrm{sym}}(h)&=\frac{x^3+y^3+z^3}{3}+u^3+v^3
                       -\frac{2h^2}{3}xyz.
\end{align*}
両式の差は正確に$h^2xyz$である。
ここで$J(F_0)=\CC[x,y,z,u,v]/(x^2,y^2,z^2,u^2,v^2)$なので、
差の非零性は平方自由単項式$xyz$だけで確認できる。
この判定はmarked変形の判定であり、
有限自己同型群で割った粗モジュライの局所座標について、同じ次数の主張を加えていない。
\end{example}

\subsection{代数そのものの可換化は次元さえ保たない}

\begin{proposition}[文字通りの可換化の失敗]
\label{em:prop:abelianization}
$t=0$、$h\ne0$とする。$B_{h,0,P}$を、全ての交換子が生成する両側イデアルで割った
可換代数を$B_{h,0,P}^{\mathrm{ab}}$と書く。このとき
\[
 B_{h,0,P}^{\mathrm{ab}}
 \simeq
 \CC[x,y,z,u_1,\ldots,u_m]
 /(x^2,y^2,z^2,P(u)-2h^2xyz/3).
\]
その被約環は$\CC[u]/(P)$であり、$\operatorname{Proj}$の次元は$m-2$である。
従って、次元$m+1$の復元先cubicをこの可換化から得ることはできない。
\end{proposition}
\begin{proof}
関係式$[y,z]=hx^2$とその巡回置換に可換性を課すと、$h\ne0$から
$x^2=y^2=z^2=0$が従う。
中心方程式の三乗項は消え、表示した最後の関係式が残る。
$x,y,z$は冪零なので、被約化後に残る関係は$P=0$である。
$P$は原点以外に臨界点を持たない非零三次式であり、特に被約である。
よって被約環と次元の主張を得る。
\end{proof}

\subsection{非楕円混合型のモデルへの適用障害}

節点型と尖点型で用いる三変数三次式は
\[
 f_{\mathrm N}=y^3+z^3+xyz,\qquad
 f_{\mathrm C}=xz^2+y^3
\]
であり、全体の可換特殊ファイバーには$ux^2$という混合項がある。
この項は、一ブロック定理で許される中心変数方向への線形平行移動に吸収できない。

\begin{lemma}[混合項$ux^2$は中心シフトで除去できない]
\label{em:lem:essential-mixing}
$\epsilon=\mathrm N,\mathrm C$のいずれについても
$x^2\notin\langle\partial_xf_\epsilon,\partial_yf_\epsilon,
\partial_zf_\epsilon\rangle\subset\CC[x,y,z]_2$である。
従って$f_\epsilon+ux^2$から$u$の一次混合項を
$(x,y,z)\mapsto(x+a u,y+b u,z+c u)$によって除去できない。
\end{lemma}
\begin{proof}
点$(x,y,z)=(1,0,0)$では両方の$f_\epsilon$の三つの偏微分が全て零であり、
$x^2=1$である。よって表示した線形空間に$x^2$は属さない。
中心シフトで$f_\epsilon$の$u$の一次係数に加わるものは
$a\partial_x f_\epsilon+b\partial_y f_\epsilon+c\partial_z f_\epsilon$だけなので、結論を得る。
\end{proof}

従って非楕円二例には、単に三変数のpoint curveを取り出して
定理\ref{er:thm:main}を適用することはできない。
候補式を得るためには、全七変数の量子関係式と中心三次元から残余圏の
rotation付き乗法データを計算するか、混合項を保つ具体的な行列因子分解同値を作る必要がある。
候補cubicが滑らかであることや、有限個のHochschild次元が一致することだけでは同値は証明されない。
また中心変数を通常の偏微分で扱う生成性の議論は、非楕円混合型では
その偏微分が代数導分であることから確かめ直す必要がある。

\subsection{圏論的入力と有限計算の役割}

系\ref{em:cor:hessian-reconstruction}と命題\ref{em:prop:determinant}の圏同値部分は
定理\ref{er:thm:main}と、その証明で用いるdg window、
periodicityを保つブロック比較、中心変数に対する外部積生成、
有限群の双対同変化\cite[Theorem 4.2, Section 8]{Elagin}を入力としている。
行列式や中心係数の一致は、これらの圏論的入力の独立監査を代替しない。
命題\ref{em:prop:second-order-failure}、命題\ref{em:prop:abelianization}、
補題\ref{em:lem:essential-mixing}の証明は上記の初等的計算で完結する。
付属の\texttt{verify\_explicit\_model\_comparison.py}は、
一般の十係数三次式に対するHessian恒等式、Hesse形への特殊化、
多重線形化行列の行列式、二次係数、十次元係数行列の行列式、
Fermat例のJacobian剰余、混合項の評価を有理数上で厳密に検算する。
この検算は圏同値、一般混合型の可換復元、残余smoothnessの計算機証明ではない。
